% Part: proof-theory
% Chapter: natural-deduction
% Section: grafting

\documentclass[../../../include/open-logic-section]{subfiles}

\begin{document}

\olfileid{pt}{nat}{gra}

\olsection{证明的嫁接}

在 \Log{N1c} 和~\Log{N1i} 中，推导是由假设~$!B$ 推出公式~$!A$ 的推导。
假设 $!B$ 本身也有一个推导。
我们如何将 $!B$ 的推导与由~$!B$ 推出 $!A$ 的推导结合起来，
从而得到一个不使用假设~$!B$ 的 $!A$ 的推导呢？

一种方案是，通过应用~\Intro{\lif} 解除假设~$!B$。
然后，我们可以将由此得到的 $!B \lif !A$ 的推导与 $!B$ 的推导结合起来，以得到
\[
\AxiomC{$\Discharge{!B}{x}$}
\DeduceC{$!A$}
\DischargeRule{\Intro{\lif}}{x}
\UnaryInfC{$!B \lif !C$}
\AxiomC{}
\DeduceC{$!B$}
\RightLabel{\Elim{\lif}}
\BinaryInfC{$!A$}
\DisplayProof
\]
这样做固然直接，但多少有点难看：为什么刚引入 $\lif$ 又立刻消去它呢？
应该可以避免这种经由 $!B \lif !A$ 的“迂回”。（这种迂回总是可以避免的，这实际上就是正规化定理的内容。）

在 \Log{N1c} 和~\Log{N1i} 中，我们实际上在此情形下可以非常轻松地避免这种迂回：
只需简单地将所有假设~$!B^x$ 替换为 $!B$ 的整个推导即可。
这种将一个推导“嫁接”到其他推导的假设上的操作非常有用，因此我们将其记录为一条定义。

\begin{defn}
若 $\delta_1$ 是一个以开放假设~$!B^x$ 推出~$!A$ 的 \Log{N1c}-推导（或 \Log{N1i}-推导），且 $\delta_2$ 是一个~$!B$ 的 \Log{N1c}-推导（\Log{N1i}-推导），则 $\Subst{\delta_1}{\delta_2}{!B^x}$ 是将 $\delta_1$ 中每一个未解除的 $!B^x$ 的出现替换为 $\delta_2$ 所得到的结果。
\end{defn}

只要 $\delta_1$ 和 $\delta_2$ 没有带相同标签的不同假设公式，且 $\delta_2$ 没有开放假设包含 $\delta_1$ 的本征变元，那么得到的 $\Subst{\delta_1}{\delta_2}{!B^x}$ 就是一个正确的推导，并且其开放假设是 $\delta_1$ 与 $\delta_2$ 的开放假设之并。在进行推导嫁接时，这些条件通常都能满足。当 $\delta_1$ 和 $\delta_2$ 是某个更大推导的子证明时，第一个条件总是满足的。如果第二个条件不满足，我们可以对 $\delta_1$ 的本征变元进行重命名，使其得以满足。

类似的定义适用于 \Log{N2c} 和~\Log{N2i}。

\begin{defn}
若 $\delta_1$ 是一个~$x:!B, \Gamma_1 \Sequent !A$ 的 \Log{N2c}-推导（或 \Log{N2i}-推导），且 $\delta_2$ 是一个~$\Gamma_2 \Sequent !B$ 的 \Log{N1c}-推导（\Log{N1i}-推导），则 $\Subst{\delta_1}{\delta_2}{x:!B}$ 是将 $\delta_1$ 中每一个公理 $x:!B \Sequent !B$ 替换为 $\delta_2$，并将 $\Gamma_2$ 添加到这些公理下方所有相继式的语境中所得到的结果。
\end{defn}

\begin{prop}
若 $\delta_1 \Proves x:!B, \Gamma_1 \Sequent
!A$ 且 $\delta_2 \Proves \Gamma_2 \Sequent !B$，则 $\Subst{\delta_1}{\delta_2}{x:!B} \Proves \Gamma_1, \Gamma_2 \Sequent
!A$，只要 $\delta_1$ 的本征变元不出现在 $\Gamma_2$ 中。
\end{prop}


\begin{proof}
  对 $\pheight{\delta}$ 施归纳。
\end{proof}

\end{document}